%============================================================
%  MIDTERM exam sheet template (announced-question-types style)
%  Use when the teacher announces the TYPES of questions (e.g. "Q1 formulate an LP,
%  Q2 prove convexity, ...") instead of a question pool, and some are proofs.
%  For question-pool exams use final_template.tex.
%  Start one with:  scripts/new_note.sh <Course> midterm "Course Name"
%  The preamble (everything above the document body) is identical in both templates.
%  Build:  tectonic <file>.tex
%
%  Layout:
%    Playbook  - one part per announced question type:
%                origbox (the announcement) -> what a full-mark answer has ->
%                guidebox (recipe) -> proofbox (proof types) -> worked example ->
%                examsolbox -> trapbox
%    Practice  - review / exercise questions in sheet order, each tagged "trains Qn",
%                same 4-part layout as the final template
%    Theory    - by lecture, tagged with the exam question it feeds; defbox holds the
%                lecturer's exact definitions
%    Appendix  - symbol dictionary, derivative / formula tables, slide index
%  Boxes: origbox, guidebox, examsolbox, proofbox, defbox, resultbox, qbox, algobox,
%         trapbox; badges \step{n}{Name} and \iterhead{k}{...}
%============================================================
\documentclass[12pt,a4paper]{article}
\usepackage[top=1.4cm,bottom=1.6cm,left=1.6cm,right=1.6cm,headsep=5pt,footskip=20pt]{geometry}
% >>> SHARED PREAMBLE (templates/_exam_preamble.tex, do not edit here)
% Shared preamble for every EXAM template (midterm, final, theory, question bank,
% cheat sheet). Edit templates/_exam_preamble.tex, then run: python3 scripts/sync_preamble.py
% That stamps it between the ">>> SHARED PREAMBLE" markers of each exam template.
% \documentclass, geometry and the "who and what" block stay per template.
\usepackage{setspace}
\usepackage{amsmath,amssymb}
\usepackage{xcolor}
\usepackage{titlesec}
\usepackage{enumitem}
\usepackage{booktabs}
\usepackage{array}
\usepackage{multicol}
\usepackage{needspace}
\usepackage{silence}
\WarningFilter{LaTeX}{Command \showhyphens has changed}
\usepackage{microtype}
\usepackage{fancyhdr}
% Fonts: tectonic runs XeTeX, so use fontspec (OpenType Latin Modern covers
% accented names like "Trần Bảo Thành"). pdflatex falls back to the classic setup.
\usepackage{iftex}
\ifTUTeX
  \usepackage{fontspec}
\else
  \usepackage[T1]{fontenc}
  \usepackage[utf8]{inputenc}
  \usepackage{lmodern}
\fi
\usepackage{parskip}
\usepackage{listings}
\usepackage{graphicx}
\usepackage{tikz}
\usetikzlibrary{arrows.meta,positioning,calc,fit,shapes.geometric,backgrounds}
% bookmarks=false: math in a heading + PDF bookmarks hangs the compile forever
\usepackage[hidelinks,bookmarks=false]{hyperref}
\usepackage{tcolorbox}
\tcbuselibrary{skins,breakable}

% ---- spacing and page breaks ----
\setstretch{1.05}
\setlength{\parskip}{5pt}
\setlength{\parindent}{0pt}
\setlist[itemize]{itemsep=2pt,topsep=3pt,parsep=0pt}
\setlist[enumerate]{itemsep=2pt,topsep=3pt,parsep=0pt}
\raggedbottom
\widowpenalty=10000 \clubpenalty=10000
\displaywidowpenalty=10000 \predisplaypenalty=10000
\setlength{\emergencystretch}{2em}

% ---- section formatting (a heading never sits alone at the bottom of a page) ----
\titleformat{\section}{\large\bfseries}{}{0pt}{\needspace{6\baselineskip}}[\titlerule]
\titleformat{\subsection}{\normalsize\bfseries}{}{0pt}{\needspace{5\baselineskip}}
\titleformat{\subsubsection}{\normalsize\itshape\bfseries}{}{0pt}{\needspace{4\baselineskip}}
\titlespacing{\section}{0pt}{14pt}{6pt}
\titlespacing{\subsection}{0pt}{10pt}{4pt}
\titlespacing{\subsubsection}{0pt}{6pt}{2pt}

% ---- part banner: \examPart{Title}{subtitle} ----
\newcommand{\examPart}[2]{\clearpage\begin{center}{\Large\bfseries #1}\\[3pt]{\normalsize #2}\end{center}\vspace{2pt}\hrule\vspace{6pt}}

% ---- math shortcuts (add your own) ----
\newcommand{\R}{\mathbb R}
\newcommand{\sign}{\operatorname{sign}}
\newcommand{\diag}{\operatorname{diag}}

% ---- small inline helpers ----
% \sref{L1 p46}            inline source tag, lighter than a box title
% \gl{term:dot}            superscript link back to a glossary item with \label{term:dot}
% \qtag{LIKELY}            tag badge: LIKELY, PAST EXAM A, PRACTICE, COMPUTE, TWIST, ...
% \exam{Midterm + Final}   which exam a section is for (put it in the heading)
% \lbl{What it is:}        bold run-in label for fixed-pattern blocks
% \examcue{"quote"}        what the lecturer stressed, verbatim
\newcommand{\sref}[1]{\textcolor{blue!50!black}{\small[#1]}}
\newcommand{\gl}[1]{\textsuperscript{\textcolor{blue!50!black}{[\ref{#1}]}}}
\newcommand{\qtag}[1]{{\setlength{\fboxsep}{1.5pt}\colorbox{black!10}{\scriptsize\bfseries\MakeUppercase{#1}}}}
\newcommand{\exam}[1]{\ \textnormal{\small(#1)}}
\newcommand{\lbl}[1]{\par\noindent\textbf{#1}\ }
\newcommand{\examcue}[1]{\par\noindent\textcolor{red!55!black}{\textbf{Exam cue:}}\ \emph{#1}\par}

% ---- plain boxes: white background, coloured streak on the left ----
\tcbset{
  streak/.style={
    enhanced, breakable, colback=white, frame hidden, boxrule=0pt, arc=0pt,
    left=8pt, right=3pt, top=3pt, bottom=3pt,
    borderline west={3pt}{0pt}{#1},
    before skip=8pt, after skip=8pt
  },
  streaktitle/.style 2 args={
    streak=#1, title={#2}, attach title to upper, after title={\par\smallskip},
    fonttitle=\small\bfseries\color{#1}
  },
  titled/.style 2 args={
    enhanced, colback=#1!2, colframe=#1!45!black, boxrule=0.6pt, arc=2pt,
    left=7pt, right=7pt, top=5pt, bottom=5pt, before skip=10pt, after skip=10pt,
    colbacktitle=#1!45!black, coltitle=white, fonttitle=\bfseries, title={#2}
  },
  bw/.style={
    enhanced, breakable, colback=white, colframe=black, boxrule=0.4pt, arc=2pt,
    left=4pt, right=4pt, top=2pt, bottom=2pt, before skip=6pt, after skip=6pt
  }
}
\newtcolorbox{qbox}{streak=black!70}
\newtcolorbox{resultbox}{streak=black!45}
\newtcolorbox{trapbox}{streaktitle={orange!80!black}{Traps}}
\newtcolorbox{algobox}[1]{streaktitle={violet!70!black}{#1}}
% defbox: the lecturer's exact wording for a definition / named list. \begin{defbox}{Lecture 2, slide 7}
\newtcolorbox{defbox}[1]{streaktitle={teal!60!black}{Lecturer's definition \textnormal{\small(#1)}}}
% proofbox: a proof written the way it should appear on the paper. \begin{proofbox}{Title}
\newtcolorbox{proofbox}[1]{titled={orange}{Proof: #1\hfill{\normalfont\small copy this structure}}, breakable}

% ---- question box: \begin{origbox}{Title}{Source, e.g. Pool \S2.1, p.~1} ----
% Not breakable on purpose: a question should never split across pages.
\newtcolorbox{origbox}[2]{titled={blue}{#1\hfill{\normalfont\small #2}}}

% ---- guideline box: the recipe, written with letters so it works for any numbers ----
\newtcolorbox{guidebox}{titled={red}{Step-by-step guideline\hfill{\normalfont\small works for any numbers}}, breakable}

% ---- answer box: \begin{examsolbox}{Question 2.1} ----
\newtcolorbox{examsolbox}[1]{titled={green}{Answer: #1\hfill{\normalfont\small what to write on the exam paper}}, breakable}

% ---- theory-course boxes ----
% notebox     black frame, print friendly (overview, decoders, case studies)
% memobox     "Memorise this" ordered list / chain, with an optional title
% scenariobox a situational / essay question with its answer skeleton
% summarybox  the last-minute checklist at the end of every sheet
% pitfallbox  exam-day pitfalls, gathered in one numbered list
\newtcolorbox{notebox}{bw}
\newtcolorbox{memobox}[1][Memorise this]{streaktitle={teal!60!black}{#1}}
\newtcolorbox{scenariobox}[1]{titled={violet}{#1\hfill{\normalfont\small situational question}}, breakable}
\newtcolorbox{summarybox}{titled={black}{Quick-reference summary\hfill{\normalfont\small last-minute checklist}}, breakable}
\newtcolorbox{pitfallbox}{titled={orange}{Exam-day pitfalls\hfill{\normalfont\small read before you start}}, breakable}

% ---- question bank: \Q{Q12}{\qtag{past exam A} Birthday attack on SHA-2} ----
% then \begin{qbankbox} question only \end{qbankbox}, then Answer / Why / Trap / Also asked as
\newcommand{\Q}[2]{\needspace{6\baselineskip}\subsection*{#1\quad #2}}
\newtcolorbox{qbankbox}{streak=blue!50!black}

% ---- step and iteration badges ----
% \step{1}{Lagrangian}   -> red "STEP 1" badge + bold name on its own line
% \iterhead{2}{at $(0.8,0)$:} -> blue "ITERATION 2" badge
\newcommand{\step}[2]{\par\medskip\noindent{\setlength{\fboxsep}{2pt}\colorbox{red!55!black}{\textcolor{white}{\small\bfseries\,STEP #1\,}}}\ \ {\bfseries #2}\par\nobreak\smallskip}
\newcommand{\iterhead}[2]{\par\medskip\noindent{\setlength{\fboxsep}{2pt}\colorbox{blue!45!black}{\textcolor{white}{\small\bfseries\,ITERATION #1\,}}}\ \ #2\par\nobreak\smallskip}

% ---- code listings (black and white, prints well) ----
\lstdefinestyle{pyStyle}{language=Python, basicstyle=\ttfamily\small, keywordstyle=\bfseries,
  commentstyle=\itshape, stringstyle=\ttfamily, breaklines=true, frame=single, numbers=left,
  numberstyle=\tiny, xleftmargin=15pt, captionpos=b, columns=fullflexible, keepspaces=true}
\lstdefinestyle{cppStyle}{style=pyStyle, language=C++}

% ---- TikZ styles for redrawing exam diagrams (block ciphers, pipelines, flows) ----
\tikzset{
  blk/.style={draw, minimum height=6.5mm, minimum width=11mm, font=\scriptsize, align=center, inner sep=2pt},
  enc/.style={draw, thick, minimum height=9mm, minimum width=18mm, font=\scriptsize\bfseries, align=center},
  xr/.style={draw, circle, minimum size=4.5mm, inner sep=0pt, font=\footnotesize},
  kb/.style={draw, fill=black!10, minimum height=7mm, minimum width=15mm, font=\scriptsize, align=center, inner sep=2pt},
  lb/.style={font=\scriptsize, align=center},
  ar/.style={-{Stealth[length=1.6mm]}}
}

% ---- university logo: \UniLogo[height] ----
% Looks for logo.png next to the document, then assets/logo.png at the repo root from
% any folder depth (templates/, <Course>/notes/, <Course>/paper/<name>/, ...).
% No logo file -> prints nothing. Swap assets/logo.png to change university.
\def\LogoFile{}
\IfFileExists{logo.png}{\def\LogoFile{logo.png}}{%
\IfFileExists{assets/logo.png}{\def\LogoFile{assets/logo.png}}{%
\IfFileExists{../assets/logo.png}{\def\LogoFile{../assets/logo.png}}{%
\IfFileExists{../../assets/logo.png}{\def\LogoFile{../../assets/logo.png}}{%
\IfFileExists{../../../assets/logo.png}{\def\LogoFile{../../../assets/logo.png}}{}}}}}
\newcommand{\UniLogo}[1][1.6cm]{\ifx\LogoFile\empty\else\includegraphics[height=#1]{\LogoFile}\fi}
% \UniLogoBreak[h]: the logo plus a line break, both left out when there is no logo file
\newcommand{\UniLogoBreak}[1][1.6cm]{\ifx\LogoFile\empty\else\includegraphics[height=#1]{\LogoFile}\\[6pt]\fi}

\pagestyle{fancy}
\fancyhf{}
\cfoot{\small\thepage}
\renewcommand{\headrulewidth}{0pt}
% <<< SHARED PREAMBLE

% ---- who and what (edit these) ----
\newcommand{\StudentName}{Your Name}
\newcommand{\StudentID}{YOUR-ID}
\newcommand{\CourseName}{Course Name}
\newcommand{\ExamName}{Midterm Exam Sheet}
\newcommand{\ExamInfo}{Open-book reference}

\begin{document}

\begin{center}
  \UniLogoBreak[1.6cm]
  {\LARGE\bfseries \CourseName: \ExamName}\\[6pt]
  {\large \ExamInfo}\\[4pt]
  {\normalsize Student: \StudentName\ (\StudentID)}
\end{center}
\hrule
\vspace{6pt}

\section{Exam playbook (read this first)}

\textbf{Exam conditions.} Open book? Printed material only? Calculator? Solver? (If no
solver: every answer below is done by hand.)

\textbf{The announced questions} (copy the teacher's message word for word):
\begin{center}
\renewcommand{\arraystretch}{1.2}
\begin{tabular}{@{}p{0.07\linewidth}p{0.43\linewidth}p{0.16\linewidth}p{0.24\linewidth}@{}}
\toprule
Q & What she will ask & Kind & Go to\\
\midrule
Q1 & e.g.\ formulate an LP, its dual, complementary slackness & computation & Playbook Q1, p.~\pageref{pbq1}\\
Q2 & e.g.\ prove a function / problem is convex & \textbf{proof} & Playbook Q2, p.~\pageref{pbq2}\\
\bottomrule
\end{tabular}
\end{center}

\textbf{Habits that save marks.}
\begin{itemize}
  \item Name the method before using it (``by the second-order condition \dots'').
  \item Write the general formula first, then plug in the numbers.
  \item Check the answer: constraints hold, multipliers have the right sign, primal value $=$ dual value.
  \item For proofs: state the claim, the method, each step with its reason, then $\square$.
\end{itemize}

\begin{resultbox}
\textbf{Hand-calculation helpers.} Small formulas you keep looking up.
\end{resultbox}

%############################################################
\examPart{Playbook Q1: name of the question type}{Computation. Lectures N--M.}
\label{pbq1}

\begin{origbox}{Q1 \ What the teacher announced}{Announcement, <date>}
Copy the exact wording of this announced question.
\end{origbox}

\textbf{What a full-mark answer contains.} A short checklist of the parts the marker looks for.

\begin{guidebox}
\step{1}{First thing to write} Written with letters, so it works for any problem.
\step{2}{Next thing} \dots
\step{3}{Check} How to verify the answer by hand.
\end{guidebox}

\subsection{Worked example}
\begin{qbox}
A typical problem of this type (from the review sheet or a past paper).
\end{qbox}
\step{1}{First thing to write} The real numbers, every line.
\step{2}{Next thing} \dots

\begin{examsolbox}{Q1 (example)}
\begin{enumerate}[label=\textbf{(\alph*)},leftmargin=2em,itemsep=6pt]
  \item Short working, final answer boxed: $\boxed{x^\star=(\dots)}$
\end{enumerate}
\end{examsolbox}

\begin{trapbox}
The mistake most people make on this type, and a quick check that catches it.
\end{trapbox}

%############################################################
\examPart{Playbook Q2: name of the proof question}{Proof. Lectures N--M.}
\label{pbq2}

\begin{origbox}{Q2 \ What the teacher announced}{Announcement, <date>}
e.g.\ ``Prove that the problem \dots\ is a convex (concave) optimization problem.''
\end{origbox}

\textbf{Which method to pick.}
\begin{center}
\begin{tabular}{@{}p{0.38\linewidth}p{0.52\linewidth}@{}}
\toprule
If you see \dots & use \dots\\
\midrule
a twice-differentiable function & second-order condition: Hessian $\succeq0$\\
a sum / max / composition of known pieces & the rules that keep convexity\\
\bottomrule
\end{tabular}
\end{center}

\begin{guidebox}
\step{1}{State the claim} What exactly is being proved, and on which domain.
\step{2}{Pick and name the method} \dots
\step{3}{Prove each piece} One line per piece, with its reason.
\step{4}{Conclude} Tie the pieces together and finish with $\square$.
\end{guidebox}

\begin{proofbox}{template for this question type}
\textbf{Claim.} \dots\par
\textbf{Method.} \dots\par
\textbf{Step 1.} \dots\ \emph{(reason)}\par
\textbf{Step 2.} \dots\ \emph{(reason)}\par
\textbf{Conclusion.} Hence \dots\ is convex. $\square$
\end{proofbox}

\subsection{Worked example}
\begin{proofbox}{a real example, filled in}
\textbf{Claim.} $f(x)=x\log x$ is convex on $x>0$.\par
\textbf{Method.} Second-order condition.\par
\textbf{Step 1.} $f'(x)=\log x+1$, \ $f''(x)=1/x$.\par
\textbf{Step 2.} For $x>0$, $f''(x)=1/x>0$ \emph{(a positive number over a positive number)}.\par
\textbf{Conclusion.} $f''>0$ on the convex domain $x>0$, so $f$ is strictly convex. $\square$
\end{proofbox}

%############################################################
\examPart{Practice questions}{Review sheet / exercises in their original order. Each says which exam question it trains.}

\section{Question 1: short name \ \textnormal{\small(trains Q1)}}

\begin{origbox}{Question 1}{Review sheet, p.~1}
Copy the question exactly as printed.
\end{origbox}

\begin{guidebox}
\step{1}{Same recipe as Playbook Q1} Point back to it, and note anything different.
\end{guidebox}

\step{1}{\dots} Worked with the real numbers.

\begin{examsolbox}{Question 1}
\begin{enumerate}[label=\textbf{(\alph*)},leftmargin=2em]
  \item \dots
\end{enumerate}
\end{examsolbox}

%############################################################
\examPart{Theory reference}{By lecture. Each section says which exam question it feeds.}

\section{Symbol dictionary (read this first)}
\begin{center}
\begin{tabular}{@{}ll@{}}
\toprule
Symbol & Read it as \dots\\
\midrule
$x^\star$ & the optimal point\\
$\nabla f$ & the gradient: the list of slopes\\
\bottomrule
\end{tabular}
\end{center}

\section{Topic name \ \textnormal{\small(Exam Q1, Lecture 1)}}

\begin{defbox}{Lecture 1, slide 5}
The lecturer's exact definition. If the slide has a list of N items, write all N.
\end{defbox}
In plain words: one or two sentences. Then a small example with numbers.

%############################################################
\examPart{Appendix}{Tables to look things up fast.}

\section{Derivative and formula tables}
\begin{center}
\begin{tabular}{@{}ll@{\qquad}ll@{}}
\toprule
$f(x)$ & $f'(x)$ & $f(x)$ & $\nabla f(x)$\\
\midrule
$x^n$ & $nx^{n-1}$ & $a^Tx$ & $a$\\
$\ln x$ & $1/x$ & $x^TAx$ & $(A+A^T)x$\\
$e^{ax}$ & $ae^{ax}$ & $\|x\|_2^2$ & $2x$\\
\bottomrule
\end{tabular}
\end{center}

\section{Slide index}
Topic $\to$ lecture and slide number, so you can find the original fast during the exam.

\end{document}
